\section{Threats to Validity}
\label{sec:threats-to-validity}

The largest limitation to validity is that allocator versions were not all benchmarked on the same machine. Data were collected from Apple M3/M5 systems, several Intel systems, and Ryzen systems. The operating systems also varied, including WSL, native Linux, and macOS. CPU architectures differed as well, with some runs using ARM64 rather than the more typical x86-64. This limitation is mitigated largely by focusing on comparisons within the same environment and machine, so most experiments emphasize CMA versus regular new under otherwise matching conditions. However, when analysis crosses machine or environment boundaries, these differences cannot be fully controlled.

Another limitation is that different teams used different matrix sizes, repetition counts, linked-list sizes, and other workload parameters. Targeting specified workload-duration categories such as `short,'' `medium,'' and ``long'' also reintroduced machine- and environment-dependent variation, since achieving a similar runtime on different systems could require different workload parameters. This issue is again mitigated by keeping most analysis at the per-machine and per-environment level, but it strongly limits the ability to generalize the results into a broad characterization of expected CMA performance. Unequal numbers of repeated runs also make statistical comparisons across experiments less directly comparable and reduce statistical power in experiments with fewer observations. The typical sample size of five runs is also likely too small for assumption tests such as Shapiro--Wilk and Levene's test to provide strong evidence that their respective assumptions hold. Therefore, the more appropriate interpretation is that no violation of the tested assumption was detected, rather than that the data were conclusively shown to possess a particular distributional property.

Run-to-run observations may also not be completely independent. Repeated experiments on the same machine can inherit machine state from earlier runs, including cache state, CPU temperature, CPU frequency, memory pressure, and other operating-system state. There is also insufficient evidence that every benchmark run occurred under the same level of unrelated system activity, since the machines may have been used intermittently for other tasks. If CMA and regular-new benchmarks were repeatedly executed in a fixed order, these effects could introduce systematic bias. Even when execution order varied, uncontrolled machine state could still introduce additional noise and affect the resulting statistical analysis.

The generated executable programs may also not be fully comparable across machines because of differences in compilers, standard libraries, and toolchains, including GCC, Clang/LLVM, MSVC, MinGW, MSYS2, and related platform-specific build environments. These toolchains can differ in optimization behavior, code generation, standard-library implementation, and the underlying implementation used by regular \texttt{new}. Therefore, some measured differences between environments may reflect toolchain or runtime-library behavior in addition to allocator behavior.

The experiments do not independently measure the individual sources of runtime. Total application runtime can contain time spent on computation, allocation and deallocation, mutex contention, free-list searches, coalescing, cache hits and misses, paging, and other effects. The benchmarks primarily measure the final application-level runtime and related operating-system metrics rather than decomposing runtime into these components. As a result, explanations that attribute observed performance changes to specific allocator mechanisms should be treated as informed interpretations of the implementation and workload behavior rather than direct measurements of those mechanisms.

Finally, the results may not generalize directly to larger production or enterprise systems. The experiments were conducted on individual workstation-class machines using relatively controlled benchmark workloads. Larger systems can introduce additional factors such as higher thread counts, long-running mixed workloads, containers or virtual machines, hypervisor-level memory management, and more complex operating-system memory pressure. These layers can change allocation latency, locality, synchronization costs, paging behavior, and the relative importance of allocator overhead. Consequently, performance improvements observed in these benchmarks should be interpreted as evidence for the tested workloads and environments rather than as a guarantee of similar improvements in large-scale production deployments.